\section{Analytic results used in the applications}
\label{app:analytic-inputs}

We record the results of \cite{W_Hodge,W_HYM} with the
hypotheses and normalizations used in Sections~\ref{sec:analytic-family}
and~\ref{sec:integrated-theorem}. The statements below are
recollections of those results. Proposition~\ref{prop:analytic-reapplication}
verifies their hypotheses for the algebraically selected family;
Theorem~\ref{thm:integrated} gives the resulting applications
of the integral reconstruction.

\subsection{Domains and the Ricci-flat input}

Fix any finite graph regularity order $m\ge4$. There is a full
positive-radius ball $B_{\rm an}\subset B_{\rm ev}$ on which the
analytic conclusions stated below hold, together with positive
bundle-parameter and metric-basin radii. These radii may depend
on the fixed central member and on $m$. No uniformity as that
central member varies is asserted. Statements for an arbitrary
fixed spatial derivative order retain their own order-dependent
constants and neighbourhoods.

For $p\in B_{\rm an}^{\circ}$ let $g_p$ be the Ricci-flat
K\"ahler metric in $c_1(H)$, where $H=\OO(2,2)|_{X_p}$.
Its volume is $V_H=H^3/6=160/3$. The metric input in
\cite[Sections~7--10]{W_Hodge} consists of a fixed central
conically singular metric $g_0$, conic radius functions
$\varrho_p$, and initial K\"ahler metrics $\widetilde g_p^R$
which equal normalized Stenzel metrics on inner node regions.
For every fixed spatial derivative order $k$, positive
$\mu>\eta$ and constants $C_k$ give
\begin{equation}\label{eq:actual-metric-input}
 \limsup_{\substack{p\to0\\p\in B_{\rm an}^{\circ}}}
 \max_{j\le k}\sup_{X_p}\varrho_p^j
 |\nabla^j(g_p-\widetilde g_p^R)|_{\widetilde g_p^R}
                         \le C_k R^{\mu-\eta}.
\end{equation}
The order is to fix small $R$, let $p\to0$, and then decrease
$R$. On every fixed compact exterior $g_p\to g_0$ in all
fixed spatial derivatives. The radius on a node chart is
comparable to $(|z|^2+|\tau_i|)^{1/3}$ and is comparable to
one on the exterior. The estimates impose no ratio condition
on the nonzero node parameters. They concern spatial derivatives,
not differentiability with respect to $p$.

\subsection{The complete Hodge comparison}

Use the central harmonic mass $G$, surviving subspace $K_H$,
and orthogonal representatives $W_H$ of
Section~\ref{sec:integrated-theorem}.
On a fixed compact exterior $C_*$ let $\Pi_*$ be the central
$L^2$ projection onto distributionally closed two-forms, with
unrestricted tangential trace. The bounded exterior observer is
\begin{equation}\label{eq:hodge-observer}
 \mathcal A_Hu=\operatorname{pr}_{W_H}^G
                          [\Pi_*(u|_{C_*})].
\end{equation}
Exterior identifications are fixed before the degeneration.
This definition applies to form-domain inputs and does not
integrate an arbitrary unclosed $L^2$ form over cycles.

The finite energy and mass are $E_H,M_H$ in
\eqref{eq:hodge-model}.

\begin{recollection}[Ricci-flat Hodge comparison]\label{thm:hodge-input}
On a sufficiently small $B_{\rm an}$ there are $\gamma>0$
and $e(p)\to0$ such that the harmonic two-form space has
dimension two, with mass converging to $G|_{K_H}$, and precisely
two positive two-form eigenvalues lie below $\gamma$.
Their values are relatively $1\pm e(p)$ equivalent to the
generalized eigenvalues of $(E_H,M_H)$. The first positive
one-form eigenvalue and third positive two-form eigenvalue are
at least $\gamma$.

The positive weak two-form space $S_{H,p}$ consists of coexact
primitive real $(1,1)$ forms, and
$\mathcal A_H:S_{H,p}\to W_H$ is an isomorphism. With
$J_H=(\mathcal A_H|_{S_{H,p}})^{-1}$, its observed forms satisfy
\begin{equation}\label{eq:hodge-relative-forms}
 \begin{split}
 (1-e)M_H&\le \langle J_H\cdot,J_H\cdot\rangle_{g_p}
                                      \le(1+e)M_H,\\
 (1-e)E_H&\le \langle\dd J_H\cdot,\dd J_H\cdot\rangle_{g_p}
                                      \le(1+e)E_H.
 \end{split}
\end{equation}
These conclusions are uniform over arbitrary node-scale ratios.
The finite positive eigenvalues
\begin{equation}\label{eq:hodge-roots}
 \mu_\pm=
 \frac{9(w_f+w_g)\pm
              \sqrt{81(w_f-w_g)^2+4w_fw_g}}{128}
\end{equation}
satisfy $\mu_-/\mu_+\le4/5$. Each actual weak eigenline is
therefore identified through $\mathcal A_H$ with its finite
generalized eigenline with error tending to zero uniformly in
$w_f/w_g$.
\end{recollection}

This is the statement of \cite[Theorem~1.1 and Section~12]{W_Hodge}
in the present notation. The operator is the real Hodge Laplacian
on its closed-manifold elliptic domain. On the coexact two-form
sector its quadratic form is $\|\dd u\|_2^2$ with $H^1$ form
domain; the harmonic and exact sectors are accounted for separately.

\subsection{Nonflat bundles and their weak section spectrum}

On the same relative family put $L_p=\OO(1,-1)|_{X_p}$.
There is a locally free holomorphic rank-four family
$\mathcal E_p(w)$, $w=(w_A,w_B)\in
\mathcal W=\operatorname{Mat}_2(\CC)^{\oplus2}$ near zero,
with split member
\[
 \mathcal E_p(0)=L_p\otimes\CC^2\oplus L_p^{-1}\otimes\CC^2.
\]
It extends over the discriminant as an algebraic bundle
construction. The metric assertions here are on smooth fibres.
Writing $D_L=c_1(L_p)$, its characteristic data are
\begin{equation}\label{eq:nonflat-characteristic-data}
 D_LH^2=0,\qquad D_L^2H=-16,\qquad
 \det\mathcal E=\OO,\quad c_2(\mathcal E)=-2D_L^2,\quad
 c_2(\mathcal E)H=32,\quad c_3(\mathcal E)=0.
\end{equation}
In particular this is a nonflat bundle family.

The local deformation germ is the full two-arrow space:
$\operatorname{Ext}^0=\mathfrak{gl}_2\oplus\mathfrak{gl}_2$,
$\operatorname{Ext}^1=\mathcal W$,
$\operatorname{Ext}^2=\mathcal W^*$ and
$\operatorname{Ext}^3=(\mathfrak{gl}_2\oplus\mathfrak{gl}_2)^*$.
Its obstruction equations vanish by the ambient construction,
despite the nonzero $\operatorname{Ext}^2$ on the threefold.
All small members are Gieseker semistable. Slope polystability
is equivalent to a closed affine orbit of
$\operatorname{GL}_2(\CC)\times\operatorname{GL}_2(\CC)$.
The small germ contains no slope-stable rank-four objects.

Use the harmonic extension masses $c_\pm$, compact group $K_B$,
Hermitian space $\mathfrak h$, and finite forms and moment map
in \eqref{eq:bundle-model}.

\begin{recollection}[Nonflat bundles and their section spectrum]\label{thm:bundle-input}
There is a real $K_B$-equivariant normalization
$\chi_p(z)=z+O(|z|^2)$, $d\chi_p=I+O(|z|)$, for which
the transported corrected Chern connection at $w=\chi_p(z)$
has mean curvature exactly $\mu_{2,p}(z)$.
Thus $\mu_{2,p}(z)=0$ gives an HE representative.
These representatives cover the sufficiently small polystable
objects; this coverage uses a separate complex-orbit correction.
The map $\chi_p$ is not asserted holomorphic or complex-orbit
preserving.

At an HE member the trace-free Hermitian section operator
$\mathcal L_{p,w}=D_{p,w}^*D_{p,w}
                 =2\bar\partial_{p,w}^*\bar\partial_{p,w}$
has exactly seven low modes, zeros included. If $\nu_j$
are the generalized eigenvalues of $(E_B,M_B)$, in increasing
order with zeros, then
\[
 (1-C|z|)\nu_j\le\lambda_j(\mathcal L_{p,w})
       \le(1+C|z|)\nu_j\quad(1\le j\le7),\qquad
 \lambda_8(\mathcal L_{p,w})\ge c>0.
\]
The kernel is the constant Hermitian stabilizer in the specified
unitary transport. The full Hermitian endomorphism statement adds
the identity, with a uniformly gapped nonconstant scalar sector.
The operator domain is $H^2$ and its form domain is $H^1$ on
each fixed smooth compact fibre.
\end{recollection}

These assertions, including the deformation and stability
statements preceding this recollection, are proved in
\cite[Sections~2--7]{W_HYM}. The four factor types, in order,
are two distinct nonzero cycles, one repeated nonzero cycle,
one nonzero cycle and the two zero vertex factors, and the
split object. Their trace-free stabilizer dimensions are
$1,3,2,7$, and their positive weak counts are $6,4,5,0$.
Nonzero simple cycles give stable rank-two factors with
$c_2H=16$; zero factors are $L_p$ and $L_p^{-1}$.
All these types are included, with arbitrary ratios of positive
weak rates and independent base and bundle limits.

\subsection{Actual flow and harmonic moduli inputs}

Fix an HE member with $w=\chi_p(z)$, and let $h_0$ be its
HE metric. Use the specified unitary transport, the split
connection $D_0$, and split $L^2$ projection $\Pi$ onto
$\mathfrak h$; put $Q=I-\Pi$. For a trace-free Hermitian
logarithm define
\[
 \|u\|_X=\|u\|_\infty+\|D_0u\|_\infty.
\]
This is a spatial norm on metrics of one fixed holomorphic bundle.

\begin{recollection}[Hermitian--Yang--Mills metric flow]\label{thm:flow-input}
For the fixed finite $m\ge4$, a common nonempty ball
$\|u_0\|_X<\delta$ of smooth determinant-normalized metrics
$h(0)=h_0e^{u_0}$ has the following properties.
There are a chosen $C^m$ graph
$Z(S)=S+\psi(S)$, a complete finite field $f$, and a
$C^1$ phase $\Gamma$ depending on the full spatial datum.
Writing $r=|w|$, one has
$\psi(0)=0$, $\Pi\psi=0$,
$\|D^j\psi\|_X\le C_mr^2$ for $j\le m$ and
$\|f\|_{C^m}\le C_mr^2$.
The graph is invariant for its cutoff extension and agrees
with the actual HYM equation on an inner region.

For every datum in the stated ball,
$S_{\rm eff}(t)=\phi_t(\Gamma(u_0))$ remains in that actual
region for all $t\ge0$, and
$h_{\rm eff}(t)=h_0e^{Z(S_{\rm eff}(t))}$ is itself an actual
solution of $h^{-1}\dot h=-2\mathrm i\Lambda F_h$.
With $D_*=\|Qu_0-\psi(\Pi u_0)\|_X$,
\begin{equation}\label{eq:actual-flow-tracking}
 \begin{split}
 \|u(t)-Z(S_{\rm eff}(t))\|_X&\le CD_*e^{-ct},\\
 \left\|\log\!\left(e^{-Z(S_{\rm eff}(t))/2}
                  e^{u(t)}e^{-Z(S_{\rm eff}(t))/2}\right)\right\|_X
                                      &\le CD_*e^{-ct}.
 \end{split}
\end{equation}
The corresponding metric-order bounds are
$e^{-CD_*e^{-ct}}h_{\rm eff}\le h\le e^{CD_*e^{-ct}}h_{\rm eff}$.
The phase satisfies $\Gamma(Z(S))=S$ and, on a smaller
forward-confined basin, $\Gamma(u(t))=\phi_t\Gamma(u_0)$.

For finite preparation set $\Gamma_T=\phi_{-T}(\Pi u(T))$.
For a fixed small coupling bound $\epsilon$ and
$L_f=\|Df\|_\infty=O(r^2)<c$,
\begin{equation}\label{eq:finite-preparation}
 |\Gamma_T-\Gamma|+
 \sup_{t\ge0}\|Z(\phi_t\Gamma_T)-Z(\phi_t\Gamma)\|_X
                    \le C\epsilon D_*e^{-(c-L_f)T}.
\end{equation}
One can take $c=\sigma_0/2$, with a uniform complementary
semigroup exponent $\sigma_0>0$. For $r>0$, preparation time
$T_r=4\sigma_0^{-1}\log(1/r)$ gives error $O(r^2D_*)$
and $T_r\lambda_j\to0$ for every positive weak rate.
At $r=0$ the finite trajectory is constant.
\end{recollection}

This is \cite[Sections~9--11]{W_HYM}. The parameters $p,w$
are fixed during every flow. The result is absolute tracking
to the exact finite field, including all later times; it is
not a relative nonlinear decay expansion.

For varying holomorphic bundle parameters, put
\[
 U=\operatorname{diag}(u_1,u_2),\qquad
 z_p(u)=(c_+^{-1/2}U,c_-^{-1/2}U),\qquad
 \mathsf W=(\ZZ/2)^2\rtimes S_2.
\]
The group acts by signs and permutation.
If $\mathcal A_p(w)$ denotes the corrected
unitary connection, define the actual class map
\[
 \mathfrak F_p([u])=[\mathcal A_p(\chi_p(z_p(u)))].
\]
\begin{recollection}[Harmonic moduli metric]\label{thm:moduli-input}
For a common sufficiently small $R_0>0$,
$\mathfrak F_p:\overline B_R/\mathsf W\to\mathcal M_p(R)$
is a homeomorphism onto its actual class image for $R\le R_0$,
with smooth parametrizations of its genuine orbit-type strata.
On each stratum,
\[
 (1-C|u|)4|\dd u|^2\le\mathfrak F_p^*g_{\rm harm}
                              \le(1+C|u|)4|\dd u|^2.
\]
Let $d_{p,R}$ use continuous paths confined to $\mathcal M_p(R)$,
piecewise smooth in its finitely many strata, allowing improper
endpoint lengths. Then
\begin{equation}\label{eq:confined-metric-input}
 \sqrt{1-CR}\,d_{\rm flat}([u],[v])
 \le d_{p,R}(\mathfrak F_p[u],\mathfrak F_p[v])
 \le\sqrt{1+CR}\,d_{\rm flat}([u],[v]),
 \quad d_{\rm flat}=2\min_{\gamma\in\mathsf W}|u-\gamma v|.
\end{equation}
The open-radius locus completes to the closed-radius locus.
Its fixed local intrinsic metric has unique tangent cone
$(\CC^2,4|\dd u|^2)/\mathsf W$ at the split object.
The rescaled distance error is $O(\rho)$ on bounded balls,
uniformly in $p$, also for the fixed-locus distance $d_{p,R_0}$.
\end{recollection}

The actual-class and full-stabilizer assertions are essential
inputs here, not consequences of a formal quotient metric.
They and the stated distance theorem are proved in
\cite[Sections~12--13]{W_HYM}. Distances in a larger ambient
moduli space are outside this conclusion.

